\documentclass[11pt,a4paper]{article}

% Pakete
\usepackage[utf8]{inputenc}
\usepackage[T1]{fontenc}
\usepackage[ngerman]{babel}
\usepackage[left=2cm,right=2cm,top=2cm,bottom=2cm]{geometry}
\usepackage{helvet}
\renewcommand{\familydefault}{\sfdefault}
\usepackage{amsmath,amssymb}
\usepackage{graphicx}
\usepackage{tikz}
\usepackage{tcolorbox}
\tcbuselibrary{skins,breakable}
\usepackage{enumitem}
\usepackage{tabularx}
\usepackage{colortbl}
\usepackage{qrcode}
\usepackage{hyperref}
\hypersetup{colorlinks=true,linkcolor=blue,urlcolor=blue}

% Fachfarbe Physik
\definecolor{physik}{HTML}{2c5aa0}
\definecolor{physiklight}{HTML}{e3f2fd}
\definecolor{loesung}{HTML}{c62828}

% Box-Stile
\tcbset{
    infobox/.style={
        colback=physiklight,
        colframe=physik,
        fonttitle=\bfseries,
        title=#1,
        boxrule=1pt,
        arc=2pt
    },
    warnbox/.style={
        colback=yellow!10,
        colframe=orange!80!black,
        fonttitle=\bfseries,
        title=#1,
        boxrule=1pt,
        arc=2pt
    }
}

% Kopfzeile
\usepackage{fancyhdr}
\pagestyle{fancy}
\fancyhf{}
\lhead{\small Physik 12 GK -- Lehrerhinweise}
\chead{\small Photoeffekt Stunde 1}
\rhead{\small 07.01.2026}
\renewcommand{\headrulewidth}{0.4pt}

% Keine Einrückung
\setlength{\parindent}{0pt}
\setlength{\parskip}{6pt}

\begin{document}

% Titel
\begin{center}
{\Large\textbf{Lehrerhinweise: Photoeffekt -- Einstieg}}\\[0.3em]
{\normalsize Stunde 1 (45 min) | Mi 07.01.2026 | Hallwachs-Versuch}
\end{center}

\vspace{0.5em}

% ============================================================
% SIMULATIONEN MIT QR-CODES
% ============================================================

\section*{Simulationen}

\begin{tcolorbox}[infobox={Haupt-Simulation: PhET Photoeffekt (Deutsch)}]
\begin{minipage}{0.7\textwidth}
\textbf{URL:} \url{https://phet.colorado.edu/de/simulations/photoelectric}

\textbf{Vorteile:}
\begin{itemize}[nosep]
\item Komplett auf Deutsch
\item Verschiedene Metalle wählbar
\item Frequenz und Intensität unabhängig regelbar
\item Elektronen-Energie ablesbar
\end{itemize}

\textbf{Hinweis:} Funktioniert nicht im Internet Explorer. Chrome/Firefox/Safari empfohlen.
\end{minipage}
\hfill
\begin{minipage}{0.25\textwidth}
\centering
\qrcode[height=2.5cm]{https://phet.colorado.edu/de/simulations/photoelectric}

\vspace{0.2cm}
{\scriptsize PhET Photoeffekt}
\end{minipage}
\end{tcolorbox}

\vspace{0.5em}

\begin{tcolorbox}[infobox={Alternative 1: LEIFI Photoeffekt}]
\begin{minipage}{0.7\textwidth}
\textbf{URL:} \url{https://www.leifiphysik.de/quantenphysik/quantenobjekt-photon/downloads/photoeffekt-simulation}

\textbf{Einsatz:} Falls PhET nicht lädt oder als Ergänzung. Zeigt auch Gegenfeldmethode.
\end{minipage}
\hfill
\begin{minipage}{0.25\textwidth}
\centering
\qrcode[height=2.5cm]{https://www.leifiphysik.de/quantenphysik/quantenobjekt-photon/downloads/photoeffekt-simulation}

\vspace{0.2cm}
{\scriptsize LEIFI Photoeffekt}
\end{minipage}
\end{tcolorbox}

\vspace{0.5em}

\begin{tcolorbox}[infobox={Alternative 2: MintApps Photoeffekt}]
\begin{minipage}{0.7\textwidth}
\textbf{URL:} \url{https://www.leifiphysik.de/quantenphysik/quantenobjekt-photon/downloads/aeusserer-photoeffekt-simulation-mintapps}

\textbf{Einsatz:} Einfachere Darstellung, gut für schwächere SuS.
\end{minipage}
\hfill
\begin{minipage}{0.25\textwidth}
\centering
\qrcode[height=2.5cm]{https://www.leifiphysik.de/quantenphysik/quantenobjekt-photon/downloads/aeusserer-photoeffekt-simulation-mintapps}

\vspace{0.2cm}
{\scriptsize MintApps}
\end{minipage}
\end{tcolorbox}

\vspace{1em}

% ============================================================
% STUNDENVERLAUF
% ============================================================

\section*{Stundenverlauf}

\begin{tabularx}{\textwidth}{|c|l|X|c|}
\hline
\rowcolor{physiklight}
\textbf{Min} & \textbf{Phase} & \textbf{Aktivität} & \textbf{SF} \\
\hline
5 & Einstieg & L zeigt Bild Hallwachs-Versuch, stellt Frage: ,,Was passiert?'' & PL \\
\hline
12 & Simulation & SuS erkunden PhET: Frequenz/Intensität variieren, Beobachtungen notieren & EA/PA \\
\hline
8 & Erarbeitung & SuS bearbeiten AB Aufgabe 1--2 (Beobachtungen, Hypothesen) & EA \\
\hline
10 & Diskussion & Sammeln der Hypothesen an Tafel, Widerspruch zum Wellenmodell herausarbeiten & PL \\
\hline
5 & Sicherung & SuS füllen Merksatz (Kasten 1) aus & EA \\
\hline
5 & Ausblick & L kündigt Einstein-Erklärung an: ,,Licht besteht aus Energiepaketen'' & PL \\
\hline
\end{tabularx}

\vspace{1em}

% ============================================================
% LÖSUNGEN
% ============================================================

\section*{Lösungen zum Arbeitsblatt}

\textbf{Aufgabe 1:}
\begin{enumerate}[label=\alph)}]
\item {\color{loesung} Elektronen werden aus dem Metall herausgelöst, Strom fließt.}
\item {\color{loesung} Keine Elektronen, obwohl die Intensität hoch ist.}
\item {\color{loesung} Elektronen treten weiterhin aus, aber weniger pro Zeiteinheit.}
\item {\color{loesung} $\boxtimes$ Frequenz}
\end{enumerate}

\textbf{Aufgabe 2:}
\begin{enumerate}[label=\alph)]
\item {\color{loesung} Nach dem Wellenmodell sollte intensives rotes Licht genug Energie liefern, um Elektronen auszulösen -- tut es aber nicht. Die Frequenz (Farbe) ist entscheidend, nicht die Intensität.}
\item {\color{loesung} Mögliche Hypothesen: Licht überträgt Energie in ,,Portionen'', die von der Frequenz abhängen. Höhere Frequenz = größere Energieportionen.}
\end{enumerate}

\textbf{Kasten 1 (Lücken):}
\begin{itemize}[nosep]
\item {\color{loesung} Elektronen}
\item {\color{loesung} Licht / elektromagnetische Strahlung}
\item {\color{loesung} Frequenz}
\item {\color{loesung} Anzahl}
\end{itemize}

\textbf{Zusatzaufgabe:}
\begin{enumerate}[label=\alph)]
\item {\color{loesung} Natrium (niedrigste Grenzfrequenz)}
\item {\color{loesung} Die Grenzfrequenz hängt vom Material ab -- genauer: von der Energie, die nötig ist, um ein Elektron aus dem Metall zu lösen (Austrittsarbeit).}
\end{enumerate}

\vspace{1em}

% ============================================================
% DIFFERENZIERUNG
% ============================================================

\section*{Differenzierung}

\begin{tabularx}{\textwidth}{|c|X|}
\hline
\rowcolor{physiklight}
\textbf{Niveau} & \textbf{Hinweise} \\
\hline
★ & Aufgabe 1 a--d mit Leitfragen: ,,Schaue auf die Elektronen. Treten welche aus?'' \\
\hline
★★ & Aufgabe 1--2 selbstständig, Merksatz ausfüllen \\
\hline
★★★ & Zusatzaufgabe: Verschiedene Metalle vergleichen, Austrittsarbeit antizipieren \\
\hline
\end{tabularx}

\vspace{1em}

% ============================================================
% HÄUFIGE FEHLER
% ============================================================

\begin{tcolorbox}[warnbox={Häufige Schülerfehler}]
\begin{tabularx}{\textwidth}{lX}
\textbf{Fehler} & \textbf{Reaktion} \\
\hline
Intensität $\leftrightarrow$ Frequenz verwechselt & Analogie: Lautstärke (Intensität) vs. Tonhöhe (Frequenz) \\
\hline
,,Rotes Licht hat keine Energie'' & Korrektur: Hat Energie, aber pro Photon zu wenig \\
\hline
,,Mehr Licht = mehr Energie'' & Stimmt für Gesamtenergie, aber nicht pro ,,Portion'' \\
\hline
\end{tabularx}
\end{tcolorbox}

\vspace{1em}

% ============================================================
% TECHNISCHE HINWEISE
% ============================================================

\section*{Technische Hinweise}

\begin{itemize}[nosep]
\item \textbf{Tablets/Laptops:} PhET läuft am besten in Chrome oder Firefox
\item \textbf{WLAN:} Sicherstellen, dass alle Geräte Zugang haben
\item \textbf{Offline-Option:} PhET-Simulation kann als HTML5-Datei heruntergeladen werden
\item \textbf{Beamer:} Für Einstieg und Sicherung Simulation am Lehrerrechner zeigen
\end{itemize}

\vspace{1em}

\hrule
\vspace{0.3em}
\begin{center}
\small Stand: 05.01.2026 | Physik 12 GK
\end{center}

\end{document}
